\section{\uppercase{Discussion and Limitations}}
\label{sec:discussion}
The evaluated contract connects discrete model decisions to authenticated bytes recovered from files. Rank coding provides useful text capacity but a strong known-model detection signal. Arithmetic A1 recovers PNG payloads under the same packet budget, yet fails on low-information text trajectories. Practical success therefore depends on serialization, finite capacity and completion together, not nominal alphabet size alone.\aifootnote

Development isolates the serialization issue. Twenty image payloads under two prompts produce 40 packet-paired sequence/static comparisons. Complete-prefix filtering recovers 40/40, while singleton filtering recovers 22/40, split 8/20 for the forest prompt and 14/20 for the soup prompt. The 18 drift failures are literal saved artifacts, not repaired texts. This supports the main implementation's consistency requirement alongside published stepwise verification. It is not an impossibility result for alternative serializers or 40 independent payloads. Twenty lossless development PNG re-saves also preserve pixels and recover exactly through fresh processes. Neither check covers JPEG recompression, cropping or paraphrasing.

Context and key have distinct roles. The 96-byte row changes model probabilities but does not acquire cryptographic secrecy from being withheld. The specified mismatch weakens one score without testing adaptive context inference or arbitrary observers. Detecting a carrier also does not recover its encrypted source. These endpoints should not be collapsed into one security claim.

Twenty groups per direction and two pinned models define a focused allocation, not broad reliability or statistical power. Shared controls and repeated scores constrain independence. Exploratory comparisons are not family-wise confirmatory tests, and a bootstrap cannot reveal population overlap absent from a separated sample. A1 is one finite-precision adaptation, with a documented stagnation limitation and no observed complete text recovery. Its outcomes cannot be generalized to all arithmetic, range or distribution-preserving methods. Changing packet sizes, compression, contexts or caps would require another comparison.

Graph replay addresses repeated image-model launch overhead. Text filtering still consumes approximately 71\%, 75\% and 89\% of fixed, gated and A1 encoding time. Faster GPU inference cannot remove that CPU bottleneck or supply information absent under the frozen cap. The 5.13-fold result applies to three matched development payloads and the pinned reference/stack, not all prospective jobs or hardware.

Exact replay is established within the tested environment, not future libraries or GPU architectures. Third-party model and code terms still apply. The image port has a sale-restricting notice, and separate checkpoint redistribution permission was not established. No weights, keys or private row bytes are distributed with the manuscript. These are reproduction and dissemination constraints, not exceptions to the observed byte-equality endpoint.
